\chapter{Expanding, Factoring, and Quadratics}

For 1--12, expand. For 13--22, factor completely. For 23--30, solve for
the variable.

\section{Core Rules}

\begin{rulebox}{The products worth recognising on sight}
\[
  (a\pm b)^{2}=a^{2}\pm2ab+b^{2},\qquad
  (a-b)(a+b)=a^{2}-b^{2},
\]
\[
  a^{3}\pm b^{3}=(a\pm b)(a^{2}\mp ab+b^{2}),\qquad
  (a+b+c)^{2}=a^{2}+b^{2}+c^{2}+2ab+2ac+2bc .
\]
Read left to right they expand; read right to left they factor. The middle
term $2ab$ is the one that is forgotten: $(a+b)^{2}\neq a^{2}+b^{2}$.
\end{rulebox}

\begin{rulebox}{The binomial theorem}
\[
  (a+b)^{n}=\sum_{k=0}^{n}\binom{n}{k}a^{n-k}b^{k},
  \qquad \binom{n}{k}=\frac{n!}{k!\,(n-k)!} .
\]
The coefficients are Pascal's triangle, so $(a+b)^{4}$ uses
$1,4,6,4,1$ and $(1-x)^{5}$ uses $1,-5,10,-10,5,-1$ --- the signs alternate
because $b=-x$. The exponents of $a$ fall while those of $b$ rise, and every
term has total degree $n$: a useful check.
\end{rulebox}

\begin{rulebox}{Factoring: a fixed order of attack}
\begin{enumerate}[leftmargin=1.6em,itemsep=0.3ex]
  \item \textbf{Common factor first}, always: $4x^{4}-4=4(x^{4}-1)$.
  \item \textbf{Recognise a pattern}: difference of squares, perfect square,
        sum or difference of cubes. Apply repeatedly ---
        $x^{4}-1=(x^{2}-1)(x^{2}+1)=(x-1)(x+1)(x^{2}+1)$.
  \item \textbf{Group} when there are four terms:
        $a^{2}-2ab+b^{2}-c^{2}=(a-b)^{2}-c^{2}=(a-b-c)(a-b+c)$.
  \item \textbf{Substitute} when the exponents are in a $2\!:\!1$ ratio:
        $u=x^{2}$ turns $x^{4}-5x^{2}+4$ into a quadratic.
  \item \textbf{Split the middle term} for $ax^{2}+bx+c$ with $a\neq1$: find
        two numbers with product $ac$ and sum $b$.
\end{enumerate}
\end{rulebox}

\begin{rulebox}{Quadratics: roots and what the discriminant says}
\[
  ax^{2}+bx+c=0
  \quad\Longrightarrow\quad
  x=\frac{-b\pm\sqrt{D}}{2a},\qquad D=b^{2}-4ac .
\]
$D>0$: two real roots. $D=0$: one double root, and the parabola touches the
axis. $D<0$: no real root, and the quadratic keeps one sign everywhere ---
which is how $x^{2}+4>0$ is proved, and how a strictly convex cost function
is recognised. If the roots are $x_{1},x_{2}$ then
$x_{1}+x_{2}=-b/a$ and $x_{1}x_{2}=c/a$, a fast check on any factorisation.
\end{rulebox}

\begin{rulebox}{The zero-product property, and its one precondition}
$AB=0$ holds if and only if $A=0$ or $B=0$. \textbf{There is no analogue for
any other number.} From $A B=12$ nothing follows about $A$ and $B$
separately. Every quadratic must therefore be brought to the form
$(\cdots)(\cdots)=0$ --- with zero on the right --- before the factors are
read off.
\end{rulebox}

\section{Worked Example}

\begin{worked}{Factor $x^{4}-5x^{2}+4$, then solve $x(x-4)=12$ and $-2P^{2}+16P-30=0$}
\wstep{Step 1. Recognise the biquadratic and substitute.}
The exponents $4$ and $2$ are in a $2\!:\!1$ ratio, so $u=x^{2}$ turns the
quartic into a quadratic:
\[
  x^{4}-5x^{2}+4=u^{2}-5u+4=(u-1)(u-4).
\]

\wstep{Step 2. Substitute back --- and keep factoring.}
\[
  (x^{2}-1)(x^{2}-4)=(x-1)(x+1)(x-2)(x+2).
\]
\wtrap{Two errors live here. Stopping at $(u-1)(u-4)$ leaves the answer in
the wrong variable; stopping at $(x^{2}-1)(x^{2}-4)$ leaves it unfactored,
since each bracket is still a difference of squares. ``Factor completely''
means until no factor factors further.}

\wstep{Step 3. Now solve $x(x-4)=12$ --- but not by reading off the factors.}
The left side is factored, which tempts one to write $x=12$ or $x-4=12$.
That is invalid: the zero-product property needs a \textbf{zero} on the
right. Expand, move everything across, and factor afresh:
\[
  x^{2}-4x=12
  \;\Longrightarrow\;
  x^{2}-4x-12=0
  \;\Longrightarrow\;
  (x-6)(x+2)=0
  \;\Longrightarrow\;
  x=6\ \text{or}\ x=-2 .
\]
Check with Vieta: the roots sum to $4=-b/a$ and multiply to $-12=c/a$.
\checkmark\ (And indeed $6\cdot2=12$, $(-2)(-6)=12$, so both satisfy the
original.)

\wstep{Step 4. And $-2P^{2}+16P-30=0$: remove the leading coefficient first.}
Dividing by $-2$ is free --- zero divided by anything is still zero --- and
it makes the factorisation visible:
\[
  P^{2}-8P+15=0\;\Longrightarrow\;(P-3)(P-5)=0\;\Longrightarrow\;P=3\ \text{or}\ P=5 .
\]
\wtrap{Divide by the leading coefficient, do not merely change signs. Writing
$2P^{2}-16P+30=0$ after multiplying by $-1$ is correct but leaves the
$2$ in place; halving as well is what turns the problem into a
one-line factorisation. And when you divide, divide \emph{every} term ---
including the constant.}

\wstep{Step 5. Read the economics.}
If this were a profit condition, $P=3$ and $P=5$ are the break-even prices,
and since the parabola opens downward ($a=-2<0$) the firm makes a positive
profit strictly \emph{between} them. The sign of $a$ carries that
information; do not discard it when you divide.
\end{worked}

\section{Practice Problems}

\begin{problems}
  \item $(x-7)(x+2)$
  \item $(3a+2b)(2a-b)$
  \item $(x+y)^{4}$
  \item $(1-x)^{5}$
  \item $\displaystyle\left(x+\frac{2}{x}\right)^{2}$
  \item $(a+b+c)^{2}$
  \item $(2x-1)(4x^{2}+2x+1)$
  \item $(x+3)(x-3)(x+1)(x-1)$
  \item $(3-2r)^{3}$
  \item $(p-q+r)(p-q-r)$
  \item $(1+i)^{4}$ (binomial expansion in $i$)
  \item $(x^{2}+x+1)(x^{2}-x+1)$
  \item Factor: $4x^{2}-9$
  \item Factor: $x^{2}-12x+36$
  \item Factor: $x^{3}+8$
  \item Factor: $x^{3}-27$
  \item Factor: $a^{2}-2ab+b^{2}-c^{2}$
  \item Factor: $6x^{2}+x-12$
  \item Factor: $K^{2}L-KL^{2}$
  \item Factor: $4x^{4}-4$
  \item Factor: $x^{4}-5x^{2}+4$ (Hint: substitute $u=x^{2}$.)
  \item Factor: $P^{2}-7P+12$
  \item Solve: $x^{2}+6x+9=0$
  \item Solve: $4x^{2}-12x+9=0$
  \item Solve: $x^{2}+x-12=0$
  \item Solve: $-2P^{2}+16P-30=0$
  \item Solve: $x^{2}+2x+5=0$
  \item Solve: $3x^{2}=27$
  \item Solve: $x(x-4)=12$
  \item Solve: $Q^{2}-10Q+21=0$
\end{problems}

\section{Solutions}

\begin{solutions}
  \item $x^{2}-5x-14$
  \item $6a^{2}+ab-2b^{2}$
  \item $x^{4}+4x^{3}y+6x^{2}y^{2}+4xy^{3}+y^{4}$
  \item $1-5x+10x^{2}-10x^{3}+5x^{4}-x^{5}$
  \item $x^{2}+4+\dfrac{4}{x^{2}}$
  \item $a^{2}+b^{2}+c^{2}+2ab+2ac+2bc$
  \item $8x^{3}-1$
  \item $x^{4}-10x^{2}+9$
  \item $27-54r+36r^{2}-8r^{3}$
  \item $(p-q)^{2}-r^{2}=p^{2}-2pq+q^{2}-r^{2}$
  \item $1+4i+6i^{2}+4i^{3}+i^{4}$
  \item $x^{4}+x^{2}+1$
  \item $(2x-3)(2x+3)$
  \item $(x-6)^{2}$
  \item $(x+2)(x^{2}-2x+4)$
  \item $(x-3)(x^{2}+3x+9)$
  \item $(a-b-c)(a-b+c)$
  \item $(3x-4)(2x+3)$
  \item $KL(K-L)$
  \item $4(x-1)(x+1)(x^{2}+1)$
  \item $(x-1)(x+1)(x-2)(x+2)$
  \item $(P-3)(P-4)$
  \item $x=-3$ (double root)
  \item $x=\tfrac32$ (double root)
  \item $x=3,\,-4$
  \item $P=3,\,5$
  \item $x=-1\pm 2i$ (complex roots)
  \item $x=\pm 3$
  \item $x=6,\,-2$
  \item $Q=3,\,7$
\end{solutions}
